\section{Conditional comparison with deposition measurements}
\label{sec:comparison}
This section compares the calculations, mainly the original Q1 cases,
with the reconstructed deposition scans. The comparisons are
conditional on the reconstructed inputs and assumptions of
\cref{sec:experiments,sec:model}; they are not independent validation.

\subsection{Comparison measures}
The primary comparison uses reconstructed relative iodine yields per
1 cm bin. Let $e_b$ be the experimental percentage and let
\begin{equation}
 m_b=(100-y_{\rm KI})\frac{a_b}{\sum_{\rm full\ column}a_b}
 \label{eq:yield}
\end{equation}
be the modelled percentage, where the stated excluded KI yield is
$y_{\rm KI}=0$, 6.6 or 7.0. Over the common scanned bins $\mathcal B$,
we report the raw discrepancy
\begin{equation}
 L_{\rm raw}=\sum_{b\in\mathcal B}|m_b-e_b|
 \quad\text{(percentage points)}.
 \label{eq:raw-l1}
\end{equation}
This measure keeps the experimental digitisation sum and the model's
explicit KI exclusion. It is neither an absolute deposited-mass error
nor a percentage relative to an experimental uncertainty.

A secondary shape diagnostic normalises both profiles over
$\mathcal B$, $\widehat m_b=m_b/\sum_{\mathcal B}m_b$ and
$\widehat e_b=e_b/\sum_{\mathcal B}e_b$, and computes the overlap
\begin{equation}
 O=100\sum_{b\in\mathcal B}\min(\widehat m_b,\widehat e_b).
 \label{eq:overlap}
\end{equation}
The overlap $O$ is conditional on the common scanned range and discards
differences in capture and total yield. The experimental normalisation
still contains KI. $O$ is a descriptive shape measure, not a validation
probability or a fitted goodness-of-fit statistic.

\subsection{Q1 peaks and complete profiles}
\Cref{tab:peaks} compares the five available Q1 iodide maxima with the
printed peak temperatures. All calculated peak temperatures lie within
the stated temperature ranges. Pure PbI$_2$ lies near the
low-temperature edge of its range, and its binned maximum moves under
refinement (\cref{sec:verification}). Steel LBE places both iodide
maxima at the measured positions, whereas silica LBE places the
PbI$_2$ and BiI$_3$ maxima 5 and 4 cm downstream, respectively.
Agreement with a printed temperature is therefore not sufficient to
establish the correct spatial location.
\begin{table}[tb]
\centering
\caption{Q1 iodide maxima on the original 18,432-cell mesh, compared
with the published peaks. Positions are in cm and temperatures in
$^\circ$C. Experimental temperature ranges are those printed by Liu et
al.\ \cite{Liu2025chemical}; they are not statistical confidence
intervals for model error.}
\label{tab:peaks}
\small\begin{tabular}{llrrrrrrr}
\toprule
Case & Species & $x_{\rm exp}$ & $x_{\rm bin}$ & $x_{3}$ & $T_{\rm exp}$ & $T_{\rm bin}$ & $T_{3}$ & $T_{\rm c}$ \\
\midrule
Pure PbI$_2$, steel & \ce{PbI2} & 36 & 36 & 37 & 368$\pm$18 & 366.5 & 350.5 & 348.7 \\
LBE, steel & \ce{PbI2} & 51 & 51 & 52 & 288$\pm$18 & 288.1 & 269.2 & 274.5 \\
LBE, steel & \ce{BiI3} & 57 & 57 & 58 & 154$\pm$14 & 166.9 & 138.8 & 158.1 \\
LBE, silica & \ce{PbI2} & 51 & 56 & 56 & 306$\pm$24 & 287.8 & 287.8 & 288.8 \\
LBE, silica & \ce{BiI3} & 59 & 63 & 64 & 127$\pm$39 & 152.8 & 119.8 & 147.1 \\
\bottomrule
\end{tabular}

\end{table}

\Cref{fig:baseline} shows the complete profiles. The raw discrepancies
$L_{\rm raw}$ are 53.90, 53.45 and 181.84 percentage points for pure
PbI$_2$, steel LBE and silica LBE, and the corresponding normalised
overlaps are 72.93\%, 71.88\% and 7.20\%. The silica case therefore
has a large profile discrepancy even though both printed temperature
ranges contain its calculated maxima.
\begin{figure*}[tb]
\centering
\includegraphics[width=\textwidth]{baseline_profiles.pdf}
\caption{Original Q1 iodine profiles compared with the digitised
published scans (grey histograms) \cite{Liu2025chemical}. Solid black
lines are total modelled wall iodine; dashed lines show its two iodide
contributions. Model totals are normalised to 100\%, 93.4\% and 93.0\%
for pure PbI$_2$, steel LBE and silica LBE, respectively, to expose the
omitted KI yield; experimental bars keep their digitised values.}
\label{fig:baseline}
\end{figure*}

The Q1 steel BiI$_3$ share of deposited iodide iodine is 65.34\%, close
to the measured iodide-only share of $61.3/(61.3+32.1)=65.63\%$.
For silica the corresponding values are 28.68\% and
$27/(27+66)=29.03\%$. These agreements are substantially conditioned
by the source masses, which were inferred from the same deposits. They
cannot serve as independent validation of gas speciation or LBE
release. Similarly, the source-driven calculations capture
91.27--98.20\% of the supplied iodine on the wall by 60 s. This does
not establish the experimental three-hour capture, because the
remaining inventory, the release history and post-release cooling are
different questions.
